\scanpage{91}
% archive sheet number: 90
\[
\left\{\begin{aligned}
\sigma_0(x,y,z)&=(\unsure{1-(x+y+z)},y,z)\\
\text{pts fixes définis par}\quad&\deleted{\unsure{x=-\tfrac12(y+z-1)}}\inserted{2x+y+z-1=0}\quad\text{(si }2\text{ inv)}
\end{aligned}\right.\tag*{(12)}
\]
\margin{\deleted{\uline{NB} en car. $2$, c'est l'hyperplan à l'infini}}
\[
\left\{\begin{aligned}
\sigma_1(x,y,z)&=(y-(1+\alpha)z,x+(1+\alpha)z,z)=(y+\alpha'z,x-\alpha'z,z)\\
\text{pts fixes définis par}\quad&\left\{\begin{aligned}
\underline{x+(1+\alpha)z=y}&\qquad\text{(\unsure{id.} en }x,y,z\text{)}\\
\deleted{\unsure{x-\alpha'z=y}}&\\
x-y-\alpha'z&=0
\end{aligned}\right.
\end{aligned}\right.\tag*{(13)}
\]
\[
\left\{\begin{aligned}\sigma_2(x,y,z)&=(x,z,y)\\\text{pts fixes définis par}\quad&y=z\end{aligned}\right.\tag*{(14)}
\]
\margin{$x=\frac12\alpha',\ y=\alpha'x=\deleted{\unsure{\frac{\alpha'^2}{2}}}\inserted{\frac{\unsure{1-\alpha'|-1-\alpha'}}2},\ z=y\unsure{\ne|=}\frac12$; \deleted{$\unsure{x=\frac\alpha2,\ y=z=-\frac1{2\alpha}}$}}

\uline{Centre} $o$ de coordonnées
\[
\deleted{\unsure{\tfrac12(\tfrac12\alpha',1,1)|\tfrac12(\alpha',1,1)}}
\]
\[
\deleted{\unsure{o=(\tfrac12,-\tfrac1{2\alpha},-\tfrac1{2\alpha})=\tfrac12(-1,-\tfrac1\alpha,-\tfrac1\alpha)}}
\]
\[
o=\frac12(\alpha',1-\alpha',1-\alpha')=\frac12(-(1+\alpha),2+\alpha,2+\alpha)\tag*{(15)}
\]

\uline{NB} l'équation $\alpha^2+\alpha-1$ implique $\alpha$ inv. et \deleted{en fait $\unsure{\Rightarrow|=}$} elle s'écrit
\[
\boxed{\alpha(\alpha+1)=1}\qquad\unsure{\text{or}|\text{ou}}\tag*{(16)}
\]
\deleted{On a $\alpha-2$ \unsure{pas inv.}}

[en fait, $\alpha(\alpha+1)(\alpha-1)(\alpha+2)\inserted{(2\alpha-1)}$ \unsure{sera} inv. $\unsure{\ldots\}}$, et $(\alpha-2)$ sera inv. ssi $5$ est inv.\dots]

Ainsi, en car. $\ne2$ le centre est donné par les coordonnées (15), et en car. $2$ c'est \inserted{sans doute} le pt à l'infini dans la direction $(-1,\frac1\alpha,\frac1\alpha)$\dots\ ].
